\documentclass[11pt]{article}
\usepackage[T1]{fontenc}
\usepackage[margin=1in]{geometry}
\usepackage{amsmath,amssymb,amsthm}
\usepackage{graphicx}
\usepackage{booktabs}
\usepackage[numbers]{natbib}
\usepackage{authblk}
\usepackage[colorlinks,allcolors=blue]{hyperref}

\newtheorem{theorem}{Theorem}
\newtheorem{lemma}[theorem]{Lemma}
\theoremstyle{definition}
\newtheorem{definition}[theorem]{Definition}

\title{Title of the Paper}
\author[1]{First Author}
\author[2]{Second Author}
\affil[1]{Department, University, City, Country\\\texttt{first@example.org}}
\affil[2]{Company, City, Country\\\texttt{second@example.com}}
\date{}

\begin{document}
\maketitle

\begin{abstract}
State the problem, what you did, and the main result in a few sentences.
\end{abstract}

\noindent\textbf{Keywords:} first keyword, second keyword, third keyword

\section{Introduction}
\label{sec:intro}
Introduce the problem and why it matters, then the contributions of this
paper. Cite related work like this~\cite{he2016deep}, or several works at
once~\cite{vaswani2017attention,he2016deep}.

\section{Related Work}
Place this work among earlier ones~\cite{knuth1984,lamport1994}.

\section{Method}
\label{sec:method}
Inline maths goes between dollar signs, such as $f(x) = x^2$. Numbered
equations use the \texttt{equation} environment:
\begin{equation}
  \label{eq:loss}
  \mathcal{L}(\theta) = \frac{1}{N} \sum_{i=1}^{N} \ell\bigl(f_\theta(x_i), y_i\bigr).
\end{equation}
Refer to it as Equation~\eqref{eq:loss}.

\section{Experiments}
\label{sec:experiments}
Table~\ref{tab:results} and Figure~\ref{fig:overview} show the results.

\begin{table}[htbp]
  \centering
  \caption{Results on the test set.}
  \label{tab:results}
  \begin{tabular}{lrr}
    \toprule
    Method & Accuracy (\%) & Time (s) \\
    \midrule
    Baseline & 81.2 & 12 \\
    Ours     & \textbf{87.5} & 15 \\
    \bottomrule
  \end{tabular}
\end{table}

\begin{figure}[htbp]
  \centering
  \includegraphics[width=0.6\linewidth]{example-image}
  \caption{Overview of the method. Upload an image and use its file name here.}
  \label{fig:overview}
\end{figure}

\section{Conclusion}
Summarise the findings and what comes next.

\section*{Acknowledgements}
Thank the people and funding that supported this work.

\bibliographystyle{plainnat}
\bibliography{references}

\appendix
\section{Additional Results}
Material that supports the paper but is not needed to follow it.

\end{document}
